\section{Layer L4: dataspaces, policy and trust}
\label{sec:layer4}

An agent that reads a plant is an internal matter. An agent that reasons across
a supply chain, comparing a supplier's carbon-footprint declaration with its
own, or correlating a component's field failures with the batch data of the firm
that made it, crosses an organisational boundary, and at that boundary two new
questions become semantic rather than administrative: \emph{may I use this?} and
\emph{can I believe it?}

\subsection{IDS and the Dataspace Protocol}

The International Data Spaces reference architecture \cite{idsa2022ram}, and the
Dataspace Protocol that grew out of it, define sovereign data exchange between
organisations:
identity and attestation for participants, connectors that mediate every
transfer, catalogues that advertise datasets with DCAT metadata, and
\emph{usage policies} attached to the data itself and enforced at the receiving
connector. Eclipse Dataspace Components provides the reference implementation.

\paragraph{Agent affordance.} Usage policy is the important part, and it is
genuinely semantic. A dataset that arrives annotated with ``may be used for
quality analysis, may not be redistributed, must be deleted after 90 days'' is
carrying constraints an agent can be \emph{prevented from violating}, in the
same mechanical way SHACL prevents an out-of-range setpoint. As agents acquire
memory and autonomy this stops being a compliance detail: an agent that
summarises a partner's confidential process data into its own long-term memory
has, silently, redistributed it.

\subsection{ODRL: policy as a checkable expression}

The Open Digital Rights Language \cite{w3c2018odrl} expresses permissions,
prohibitions and duties over assets, with constraints (purpose, time, spatial
scope, count) and parties.
It is, structurally, SHACL for entitlement, and it is the only L4 technology
scoring $3$ on U6. The parallel is worth taking seriously as an architecture:
the same lift-and-validate pattern of \S\ref{sec:shacl} applies unchanged, with
the ODRL policy graph in place of the shapes graph, and the agent's proposed
\emph{use} in place of its proposed \emph{action}. We regard the unification of
action guardrails and usage guardrails under one validation step as the natural
next step for industrial agent middleware.

\subsection{Catena-X and Manufacturing-X}

Catena-X applies the dataspace idea to the automotive value chain and, crucially,
supplies versioned, machine-readable semantic content --- its aspect models ---
that makes shared data interpretable; Manufacturing-X generalises the pattern
across manufacturing sectors. (The underlying federation framework, Gaia-X,
supplies trust and compliance machinery rather than data models, and we do not
assess it separately.)

\paragraph{Assessment.} From an agent's standpoint the significant contribution
of the -X initiatives is not the infrastructure but the insistence that
\emph{every shared dataset has a versioned, machine-readable semantic model}.
That single governance rule does more for cross-company agent reasoning than any
protocol choice: it means an agent encountering a partner's data can retrieve the
aspect model, generate a typed accessor and validate what it receives, instead of
inferring structure from examples.

\subsection{Summary of L4, and a caution}

L4 is the least mature layer and the one furthest from the plant floor, and we do
not exercise it in the experiments. We include it because the trajectory is
clear: agents are moving from reading one plant to reasoning across many
organisations, and the semantics of \emph{permission}, \emph{provenance} and
\emph{trust} (the latter carried by dataset catalogues and verifiable credentials)
will become as operationally important as the semantics of pressure. The caution is
that L4's own vocabularies remain thin on U3 and U5, policy without a domain
model still cannot tell an agent whether an action is physically safe, only
whether it is contractually allowed.
